@ID: OW17D1P1I
@BIDANG: Analisis Real

Diberikan fungsi tak nol $f:D\to\mathbb R$ dan fungsi $g:D\to\mathbb R$ dengan $D\subseteq\mathbb R$ sedemikian sehingga $\dfrac{g(x)}{f(x)}\leq 1$, $\forall x\in D$, berilah contoh fungsi $f$ dan $g$ yang menunjukkan bahwa belum tentu berlaku

$$
\sup_{x\in D} g(x)\leq \inf_{x\in D} f(x).
$$

@ID: OW17D1P2I
@BIDANG: Analisis Real

$$
\lim_{n\to\infty}\left[\left(\frac1n\right)^n+\left(\frac2n\right)^n+\cdots+\left(\frac nn\right)^n\right]=\cdots
$$

@ID: OW17D1P3I
@BIDANG: Analisis Real

Jika barisan bilangan real positif $(a_n)$ memenuhi $\displaystyle\sum_{n=1}^{\infty} a_n<\infty$, maka

$$
\lim_{n\to\infty} n^{-1}\left(\sum_{k=1}^n \sqrt{a_k}\right)^2=\cdots
$$

@ID: OW17D1P4I
@BIDANG: Analisis Real

Deret fungsi $\displaystyle\sum_{k=1}^{\infty}\frac{(-x)^{k+1}}{k}$ konvergen ke $\ldots$ dengan interval konvergensi $\ldots$

@ID: OW17D1P5I
@BIDANG: Analisis Real

Diketahui fungsi $f$ terintegral pada $[a,b]$ dan $S=\{x\in(a,b):f \text{ kontinu di } x\}$. Jika $p\in S$ dan $\displaystyle\int_a^b f^2(x)\,dx=0$, maka $p^2+f^2(p)=\cdots$

@ID: OW17D1P6I
@BIDANG: Analisis Real

Diketahui fungsi $f:[0,1]\to\mathbb R$ kontinu dan $f'(x)$ ada untuk setiap $x\in(0,1)$. Jika untuk setiap $x\in(0,1)$ berlaku $0\leq f'(x)\leq 2f(x)$, maka rumus $f$ pada $[0,1]$ adalah $f(x)=\ldots$

@ID: OW17D1P7I
@BIDANG: Analisis Real

Diketahui fungsi $f:[a,b]\to[a,b]$ kontinu dan naik. Jika $f(a)=a$ dan $E=\{x\in[a,b]:f(x)\geq x\}$, maka $f(E)=\ldots$

@ID: OW17D1P8I
@BIDANG: Analisis Real

Diberikan fungsi $f:\mathbb R\to\mathbb R$, dengan $f(x)=1+a_1\sin x+a_2\sin 2x+\cdots+a_n\sin nx$, untuk setiap $x\in\mathbb R$. Jika $|f(x)-1|=2\sin 2x$, maka nilai $|a_1+2a_2+\cdots+na_n|=\ldots$

@ID: OW17D1P9U
@BIDANG: Analisis Real

Diketahui fungsi $f$ mempunyai turunan hingga tingkat ke-$2$ pada $[0,1]$ dan fungsi $g$ didefinisikan dengan $g(x)=f(x)+f(1-x)$. Jika $f''(x)>0$, untuk setiap $x\in[0,1]$, buktikan bahwa $g$ turun pada $\left[0,\dfrac12\right]$.

@ID: OW17D1P10U
@BIDANG: Analisis Real

Diketahui fungsi $f:[0,1]\to\mathbb R$ terdiferensial dan tidak ada $x$ dengan $f(x)=f'(x)=0$. Tunjukkan bahwa himpunan $\{x\in[0,1]:f(x)=0\}$ berhingga.

@ID: OW17D1P11U
@BIDANG: Analisis Real

Diketahui fungsi $f:[0,1]\to\mathbb R$ mempunyai turunan hingga tingkat ke-$2$ pada $(0,1)$. Jika $f(0)=f(1)=0$ dan $f''+2f'+f\geq 0$ pada $(0,1)$, buktikan bahwa $f(x)\leq 0$, untuk setiap $x\in[0,1]$.

@ID: OW17D1P12I
@BIDANG: Kombinatorika

Pada sebuah pesta pernikahan terdapat enam orang (termasuk pengantin) yang hendak berfoto. Banyak cara menata pose foto dalam satu baris dari keenam orang tersebut sedemikian sehingga pengantin berdiri tidak saling berdekatan adalah $\ldots$

@ID: OW17D1P13I
@BIDANG: Kombinatorika

Sebuah rangkaian digit biner adalah sebuah barisan yang terdiri dari $1$ dan $0$. Banyaknya rangkaian digit biner yang terdiri atas tepat delapan digit $0$ dan tepat sepuluh digit $1$ sedemikian sehingga setiap kemunculan digit $0$ segera diikuti oleh digit $1$ adalah $\ldots$

@ID: OW17D1P14I
@BIDANG: Kombinatorika

Pada sebuah lemari terdapat $25$ helai baju yang terdiri atas $4$ ukuran. Lima helai baju berukuran S, $4$ helai baju berukuran M, $9$ helai baju berukuran L, dan $7$ helai baju berukuran XL. Untuk menjamin telah terambil $7$ helai baju berukuran sama, maka sedikitnya total helai baju yang harus diambil dari lemari adalah $\ldots$

@ID: OW17D1P15I
@BIDANG: Kombinatorika

Sebuah keluarga besar beranggotakan $14$ orang anak yang terdiri dari dua kelahiran kembar tiga identik, tiga kelahiran kembar dua identik, dan dua anak yang lain. Bila kembar identik tak dapat dibedakan, maka banyak pose foto berdiri dalam satu baris dari $14$ orang anak tersebut adalah $\ldots$

@ID: OW17D1P16I
@BIDANG: Kombinatorika

Solusi dari relasi rekurensi $a_n=2a_{n-1}+3a_{n-2}$ dengan $a_0=1$ dan $a_1=2$ adalah $\ldots$

@ID: OW17D1P17I
@BIDANG: Kombinatorika

Banyak cara menugaskan $5$ pekerjaan berbeda ke $4$ orang pegawai berbeda sedemikian sehingga setiap pegawai ditugaskan ke paling sedikit satu pekerjaan adalah $\ldots$

@ID: OW17D1P18I
@BIDANG: Kombinatorika

Dalam bentuk yang paling sederhana fungsi pembangkit biasa (ordinary generating function), $g(x)$, dari barisan $(0,1,0,1,0,1,\ldots)$ adalah $\ldots$

@ID: OW17D1P19I
@BIDANG: Kombinatorika

Tentukan semua solusi $(a,b,c)$ dari Persamaan Diophantine $2^a+5^b=c^2$.

@ID: OW17D1P20U
@BIDANG: Kombinatorika

Seorang petinju mempunyai waktu $75$ minggu untuk mempertahankan gelar. Untuk itu pelatih menjadwalkan program latih tanding. Pelatih merencanakan sedikitnya terdapat satu latih tanding dalam satu minggu, tetapi tidak lebih dari total $125$ latih tanding dalam periode $75$ minggu. Perlihatkan ada periode waktu yang terdiri atas beberapa minggu berturutan sehingga terdapat tepat $24$ latih tanding dalam periode waktu tersebut.

@ID: OW17D1P21U
@BIDANG: Kombinatorika

Diberikan bilangan bulat $n\geq 5$. Tuliskan sebuah argumentasi kombinatorial untuk memperlihatkan bahwa

$$
\binom{2n}{5}=2\binom{n}{5}+2n\binom{n}{4}+(n^2-n)\binom{n}{3}.
$$

@ID: OW17D1P22U
@BIDANG: Kombinatorika

Suatu sisi $e$ di graf $G$ dikatakan suatu cut edge jika jumlah komponen dari $G\setminus e$ lebih dari jumlah komponen dari $G$. Buktikan bahwa, suatu sisi $e$ adalah cut edge di $G$ jika dan hanya jika $e$ tidak termuat di setiap lingkaran di $G$.

@ID: OW17D2P1I
@BIDANG: Aljabar Linear

Misalkan $K$ dan $L$ dua subruang berbeda dari ruang vektor real $V$. Jika $\dim(K)=\dim(L)=4$, maka dimensi minimal yang mungkin untuk $V$ adalah $\ldots$

@ID: OW17D2P2I
@BIDANG: Aljabar Linear

Misalkan $P_2$ adalah ruang polinom real berderajat paling tinggi $2$. Koordinat $x^2$ terhadap basis $\{x^2+x,\ x+1,\ x^2+1\}$ di $P_2$ adalah $\ldots$

@ID: OW17D2P3I
@BIDANG: Aljabar Linear

Subruang $U$ dan $W$ dari ruang vektor $\mathbb R^5$ masing-masing dibangun oleh $\{(1,3,-2,2,3),(1,4,-3,4,2),(2,3,-1,-2,9)\}$ dan $\{(1,3,0,2,1),(1,5,-6,6,3),(2,5,3,2,1)\}$. Salah satu basis dari subruang $U\cap W$ adalah $\ldots$

@ID: OW17D2P4I
@BIDANG: Aljabar Linear

Dengan hasil kali dalam $\langle A,B\rangle=\operatorname{tr}(B^tA)$, $A,B\in\mathbb R^{2\times 2}$,

$$
\left\{\begin{pmatrix}0 & 1\\ -1 & 0\end{pmatrix},\ \begin{pmatrix}1 & 1\\ a & 0\end{pmatrix}\right\}
$$

adalah himpunan ortogonal jika dan hanya jika $a=\ldots$

@ID: OW17D2P5I
@BIDANG: Aljabar Linear

Inti transformasi linier $T:\mathbb R^4\to\mathbb R^3$ dibangun oleh $\{(1,2,3,4),(0,1,1,1)\}$. Jika $T(a,b,c,d)=(a+b-c,\ x,\ 0)$, maka $x=\ldots$

@ID: OW17D2P6I
@BIDANG: Aljabar Linear

Misalkan $A=\begin{pmatrix}0 & 1\\ 0 & -1\end{pmatrix}$. Misalkan $T$ operator linier pada $\mathbb R^{2\times 2}$ dengan aturan $T(X)=AX-XA$, $\forall X\in\mathbb R^{2\times 2}$. Maka $\operatorname{rank}(T)=\ldots$

@ID: OW17D2P7I
@BIDANG: Aljabar Linear

Matriks $\begin{pmatrix}w & 1\\ -1 & 3\end{pmatrix}$ memiliki dua nilai eigen yang sama jika dan hanya jika $w\in S$. Maka $S=\ldots$

@ID: OW17D2P8I
@BIDANG: Aljabar Linear

Misalkan $V$ ruang vektor fungsi-fungsi $ae^{3x}\sin x+be^{3x}\cos x$. Transformasi $T:V\to V$ didefinisikan $T(f)=f'+f$ untuk setiap $f\in V$. Matriks representasi $T$ terhadap basis $\{e^{3x}\sin x,\ e^{3x}\cos x\}$ adalah $\ldots$

@ID: OW17D2P9U
@BIDANG: Aljabar Linear

Misalkan $A,B,C,D$ matriks-matriks berukuran $n\times n$. Misalkan pula $A$ memiliki balikan dan $AC=CA$. Buktikan bahwa

$$
\det\begin{pmatrix}A & B\\ C & D\end{pmatrix}=\det(AD-CB).
$$

@ID: OW17D2P10U
@BIDANG: Aljabar Linear

Misalkan $(v_1,v_2,\ldots,v_n)\in\mathbb R^n$ dan $A=(a_{ij})$ adalah matriks $n\times n$ dengan $a_{ij}=v_iv_j$. Jika $\operatorname{rank}(A)=1$, tentukanlah nilai $k$ yang memenuhi

$$
(I+A)^{-1}=I-\frac1k A.
$$

@ID: OW17D2P11U
@BIDANG: Aljabar Linear

Misalkan $v=\begin{pmatrix}2\\1\end{pmatrix}$, $A=vv^T$ dan

$$
K=\{x\in\mathbb R^2 \mid Ax=v\}.
$$

Tentukan $\min\{\|x\|_2 \mid x\in K\}$.

@ID: OW17D2P12I
@BIDANG: Analisis Kompleks

Berapa banyak akar berbeda dari persamaan $z^{12}=1$ yang bukan merupakan bilangan real?

@ID: OW17D2P13I
@BIDANG: Analisis Kompleks

Diketahui $f:\mathbb C\to\mathbb C$ adalah fungsi analitik dengan $f(z)=u(x)+iv(y)$ untuk setiap bilangan kompleks $z=x+iy$. Jika $f(20)=17$ dan $f(17)=20$ maka nilai dari $f(2017)$ adalah $\ldots$

@ID: OW17D2P14I
@BIDANG: Analisis Kompleks

Untuk sebarang bilangan kompleks $a$ dan bilangan real positif $r$, didefinisikan

$$
D_r^a=\{z\in\mathbb C : |z-a|<r\}.
$$

Jika fungsi $T(z)=\dfrac{z}{z+1}$ memenuhi $T^{-1}(D_r^0)=D_{2017r}^a$ maka $a=\ldots$

@ID: OW17D2P15I
@BIDANG: Analisis Kompleks

Misalkan $f$ adalah suku banyak berderajat $2$ yang memenuhi

$$
\frac{1}{2\pi i}\int_{|z|=2}\frac{zf'(z)}{f(z)}\,dz=0,\qquad \frac{1}{2\pi i}\int_{|z|=2}\frac{z^2f'(z)}{f(z)}\,dz=-2.
$$

Jika $f(0)=2017$ maka rumus eksplisit dari $f(z)$ adalah $\ldots$

@ID: OW17D2P16U
@BIDANG: Analisis Kompleks

Misalkan $\lambda$ bilangan kompleks yang memenuhi $\lambda^{2017}=1$ dan $\lambda\neq 1$.

(a) Buktikan bahwa $1,\lambda,\lambda^2,\lambda^3,\ldots,\lambda^{2016}$ semuanya berbeda.

(b) Hitung nilai dari $(1+\lambda)(1+\lambda^2)(1+\lambda^3)\cdots(1+\lambda^{2016})$.

@ID: OW17D2P17U
@BIDANG: Analisis Kompleks

Diberikan bilangan real positif $M$. Misalkan fungsi analitik $f:D\to\mathbb C$ dengan $D=\{z:|z|\leq 1\}$ mempunyai bentuk deret

$$
f(z)=\sum_{n=0}^{\infty} a_nz^n
$$

dan memenuhi $|f(z)|\leq M$ untuk setiap $z\in D$.

(a) Untuk setiap $n\geq 1$ dan $0<r<1$, buktikan bahwa

$$
|a_n|=\frac{1}{2\pi}\left|\int_{|z|=r}\frac{f(z)}{z^{n+1}}\,dz\right|.
$$

(b) Dengan menggunakan (a) dan mengambil $r\to 1$, buktikan bahwa $|a_n|\leq M$ untuk setiap $n\geq 0$.

@ID: OW17D2P18I
@BIDANG: Struktur Aljabar

Banyaknya unit di ring $\mathbb Z_{2^n}$ adalah $\ldots$

@ID: OW17D2P19I
@BIDANG: Struktur Aljabar

Misalkan $S_5$ adalah grup permutasi atas $\{1,2,3,4,5\}$. Banyaknya unsur berorde $2$ di $S_5$ adalah $\ldots$

@ID: OW17D2P20I
@BIDANG: Struktur Aljabar

Banyaknya subgrup dari $\mathbb Z_2\times\mathbb Z_4$ berorde $4$ adalah $\ldots$

@ID: OW17D2P21I
@BIDANG: Struktur Aljabar

Misalkan $\mathbb F_2$ adalah lapangan (field) dengan dua unsur. Semua polinom tereduksi berderajat $5$ di $\mathbb F_2[x]$ yang tidak memiliki akar adalah $\ldots$

@ID: OW17D2P22U
@BIDANG: Struktur Aljabar

Misalkan $H$ suatu subgrup normal hingga dari $G$. Misalkan pula $g\in G$ berorde $n$ dan unsur di $H$ yang komutatif dengan $g$ hanyalah unsur identitas $e$.

(a) Buktikan bahwa pemetaan $f:H\to H$ dengan $f(h)=g^{-1}h^{-1}gh$ merupakan suatu bijeksi.

(b) Tunjukkan bahwa semua unsur di koset $gH$ semuanya berorde $n$.

@ID: OW17D2P23U
@BIDANG: Struktur Aljabar

Buktikan bahwa $I$ merupakan ideal maksimal di gelanggang $R$ jika dan hanya jika terdapat suatu lapangan (field) $F$ dan homomorfisma ring $f:R\to F$ yang surjektif sedemikian sehingga $I=\operatorname{Ker} f$.